% =====================================================================
%  1bat-ud6-10.tex  —  SESSIÓ 10: La paradoxa del test mèdic (falsos positius)
%  (sense preàmbul: s'inclou des de main.tex)
% =====================================================================
\horatitol{Sessió 10 — La paradoxa del test mèdic (falsos positius)}%
{Un test «99\% fiable» dóna positiu. Quina probabilitat real hi ha d'estar malalt? La resposta, sorprenent, depèn de quanta gent té de debò la malaltia.}

\subsection{Idea clau}

Apliquem la probabilitat condicionada a un dels resultats més contraintuïtius i útils de
tota la matemàtica: la \textbf{paradoxa dels falsos positius}. Un tècnic de serveis a la
comunitat ha d'entendre-la per acompanyar persones que reben un resultat «positiu» en un
test (de salut, de cribratge) i evitar alarmes desproporcionades. La pregunta clau:
\emph{si un test molt fiable em dóna positiu, quina probabilitat real tinc d'estar
malalt?}

\begin{idea}
\textbf{Per què un test molt fiable pot enganyar}
\begin{itemize}[leftmargin=1.4em, itemsep=2pt]
  \item Quan una malaltia és \textbf{rara}, hi ha \emph{moltíssima} gent sana.
  \item Encara que el test falli poc amb els sans (pocs \textbf{falsos positius} per
        persona), com que els sans són tantíssims, el \textbf{nombre} de falsos positius
        pot superar el de malalts detectats.
  \item Per això $P(\text{malalt}\mid\text{positiu})$ pot ser molt \textbf{més baixa} del
        que sembla.
\end{itemize}
\textbf{L'eina:} construir una \textbf{taula de contingència} amb una població concreta
fa visible la paradoxa.
\end{idea}

\subsection{Exemples}

\begin{exemple}[la taula de contingència]
Imaginem $\num{10000}$ persones. La malaltia afecta l'\textbf{$1\%$} (prevalença), de
manera que hi ha $100$ malalts i $\num{9900}$ sans. El test detecta el $99\%$ dels
malalts (\textbf{sensibilitat}) i s'equivoca amb el $5\%$ dels sans (\textbf{falsos
positius}). Comptem:

\begin{center}
\renewcommand{\arraystretch}{1.3}
\begin{tabular}{l c c c}
\toprule
 & \textbf{Test +} & \textbf{Test $-$} & \textbf{Total} \\
\midrule
\textbf{Malalts} & $99$ & $1$ & $100$ \\
\textbf{Sans} & $495$ & $\num{9405}$ & $\num{9900}$ \\
\midrule
\textbf{Total} & $594$ & $\num{9406}$ & $\num{10000}$ \\
\bottomrule
\end{tabular}
\end{center}

Dels malalts, $99$ donen positiu ($99\%$). Dels sans, el $5\%$ ($495$) dóna positiu per
error. En total, \textbf{$594$ positius}, però \emph{només $99$ estan malalts de debò}.
\end{exemple}

\begin{exemple}[la probabilitat que sorprèn]
Si he donat positiu, quina probabilitat tinc d'estar realment malalt? És
$P(\text{malalt}\mid\text{positiu})$: dels $594$ positius, només $99$ ho estan:
\[
P(\text{malalt}\mid\text{positiu})=\frac{99}{594}\approx 0,167=16,7\%.
\]
\textbf{Sorprenent:} amb un test que «encerta el $99\%$», un positiu només significa un
$16,7\%$ de probabilitat real d'estar malalt. El motiu: hi ha tantíssima gent sana que
els falsos positius ($495$) superen de molt els malalts detectats ($99$).
\end{exemple}

\subsection{Exercicis resolts}

\begin{exresolt}[llegir la taula]
Amb la taula de l'exemple, calcula: (a) quants falsos positius hi ha; (b) la
probabilitat d'estar sa havent donat positiu.
\solucio
(a) Els \textbf{falsos positius} són sans que donen positiu: $495$.

(b) $P(\text{sa}\mid\text{positiu})$: dels $594$ positius, $495$ estan sans:
\[
P(\text{sa}\mid\text{positiu})=\frac{495}{594}\approx 0,833=83,3\%.
\]
(Coherent: $16,7\%+83,3\%=100\%$, perquè o s'està malalt o sa.)
\resposta{(a) $495$ falsos positius; (b) $83,3\%$ de probabilitat d'estar sa tot i el
positiu.}
\end{exresolt}

\begin{exresolt}[com canvia si la malaltia és més freqüent]
Repeteix el càlcul si la malaltia afecta el \textbf{$20\%$} de la població (no l'$1\%$),
amb el mateix test ($99\%$ de sensibilitat, $5\%$ de falsos positius), sobre
$\num{10000}$ persones.
\solucio
Ara hi ha $\num{2000}$ malalts i $\num{8000}$ sans. Positius:
\begin{itemize}[leftmargin=1.4em, itemsep=1pt]
  \item Malalts positius: $99\%$ de $\num{2000}=\num{1980}$.
  \item Falsos positius: $5\%$ de $\num{8000}=400$.
\end{itemize}
Total positius: $\num{1980}+400=\num{2380}$.
\[
P(\text{malalt}\mid\text{positiu})=\frac{\num{1980}}{\num{2380}}\approx 0,832=83,2\%.
\]
Amb una malaltia \textbf{freqüent}, el mateix test és molt més fiable ($83,2\%$ contra
$16,7\%$). \textbf{La prevalença ho canvia tot.}
\resposta{$P(\text{malalt}\mid\text{positiu})\approx 83,2\%$; la freqüència de la
malaltia és decisiva.}
\end{exresolt}

\subsection{Exercicis per fer}

\begin{exercicis}
\begin{exlist}
  \item D'una població de $\num{10000}$ persones amb una malaltia que afecta el $2\%$:
        quants malalts i quants sans hi ha?
  \item Amb la població de l'exercici 1 i un test que detecta el $90\%$ dels malalts i
        dóna un $10\%$ de falsos positius, completa la taula de contingència (malalts +,
        malalts $-$, sans +, sans $-$).
  \item Amb la taula de l'exercici 2, calcula el total de positius i
        $P(\text{malalt}\mid\text{positiu})$. T'esperaves un valor tan baix?
  \item Calcula també $P(\text{sa}\mid\text{positiu})$ i comprova que, sumada amb
        l'anterior, dóna $1$.
  \item \textbf{(Interpretació)} Per què, en una malaltia rara, els \textbf{falsos
        positius} poden ser molt més nombrosos que els malalts detectats, encara que el
        test falli poc \emph{per persona}?
  \item \textbf{(Raonament)} Com a tècnic, què li diries a una persona alarmada perquè
        un cribratge d'una malaltia rara li ha donat positiu? Per què és important
        \textbf{no} interpretar un positiu com una certesa i fer-ne una segona prova?
\end{exlist}
\end{exercicis}

% ---------------------------------------------------------------------
\fullsolucions{Sessió 10}
\begin{solucions}
\begin{sollist}
  \item $2\%$ de $\num{10000}=200$ malalts; $\num{9800}$ sans.
  \item Malalts +: $90\%$ de $200=180$; malalts $-$: $20$. Sans +: $10\%$ de
        $\num{9800}=980$; sans $-$: $\num{8820}$.
  \item Total positius: $180+980=\num{1160}$.
        $P(\text{malalt}\mid\text{positiu})=\dfrac{180}{\num{1160}}\approx 0,155=15,5\%$.
        És un valor sorprenentment baix, per la mateixa raó de l'exemple (molts sans).
  \item $P(\text{sa}\mid\text{positiu})=\dfrac{980}{\num{1160}}\approx 0,845=84,5\%$.
        Comprovació: $15,5\%+84,5\%=100\%$.
  \item Perquè els sans són \textbf{moltíssims} comparats amb els malalts: un percentatge
        petit de falsos positius aplicat a una població enorme de sans dóna un
        \textbf{nombre} gran de falsos positius, que pot superar de molt els pocs
        malalts (que són pocs perquè la malaltia és rara).
  \item Resposta oberta. Idea: li explicaria que un positiu en un cribratge d'una
        malaltia rara \textbf{no} vol dir que estigui malalt; sovint la probabilitat
        real és baixa pels falsos positius. Cal mantenir la calma i fer una \textbf{segona
        prova} (més específica) per confirmar, perquè una sola prova positiva no és una
        certesa.
\end{sollist}
\end{solucions}
